\documentclass[12pt]{article}
\usepackage[utf8]{inputenc}
\usepackage[T1]{fontenc}
\usepackage{hyperref}
\title{Tanema Language Resources for K-12 Curriculum Development}
\author{Compiled by Claude Sonnet 5 (claude-sonnet-5) for Andr\'as Kornai}
\date{}
\begin{document}
\maketitle

\section{Introduction}
Tanema (ISO 639-3 tnx) is an Oceanic language of Vanikoro, Temotu Province, Solomon Islands. The search results describe it as nearly extinct: a 2012 count gives a single speaker, and earlier figures from the linguist Alexandre Fran\c{c}ois (2005) give only a handful of elders. UNESCO is reported to classify it as critically endangered. This is more severe than the reference description, which says only that elderly speakers use it. Alternate names reported are Tetawo and Tetau. No school materials were found.

\section{Primary Educational Resources}
No primers, textbooks or school programmes for Tanema were found in the search results. A K-12 curriculum would have to be built from the linguistic documentation below.

\section{Linguistic and Reference Materials}
\begin{itemize}
\item Wikipedia overview of Tanema: \url{https://en.wikipedia.org/wiki/Tanema_language}
\item Fran\c{c}ois (2009), paper on the Vanikoro languages (preprint): \url{https://halshs.archives-ouvertes.fr/halshs-00525557v1/file/AlexFrancois_2009_Vanikoro-languages_Festschrift-MDR_preprint.pdf}
\item Fran\c{c}ois (2021), introduction to the Teanu dictionary on Dictionaria. It is for the neighbouring Teanu language and includes some Tanema lexical data: \url{https://marama.huma-num.fr/data/AlexFrancois_2021_Dictionaria_Teanu-dictionary_Introduction.pdf}
\item Dictionaria entry for the Teanu dictionary: \url{https://dictionaria.clld.org/contributions/teanu}
\item Fran\c{c}ois's Vanikoro notes page: \url{https://marama.huma-num.fr/AF-Vanikoro-notes_e.htm}
\end{itemize}

\section{Religious and Community Resources}
None found.

\section{Digital and Audio Resources}
A search summary says Fran\c{c}ois's 2022 ``Presentation of the Tanema language'' is held in the Pangloss Collection as open-access audio. I did not see a direct URL for it, so it is not listed as a resource.

\section{Curriculum Development Context}
With one or very few speakers, a curriculum in Tanema is not currently feasible as a living-language programme. Revitalisation would rest on the Fran\c{c}ois documentation, and on Teanu, which is now the main language of Vanikoro.
\end{document}